\section{Conclusão}
\label{sec:conclusao}

O primeiro objetivo específico era analisar se a IA reduz o tempo e o esforço de criar os \textit{templates}. Os dados confirmam isso de forma direta: a geração automatizada leva 2,7 minutos sem contexto e 3,3 minutos com contexto, contra 77,5 minutos de desenvolvimento manual líquido. A diferença entre IA e manual é muito maior do que o custo de acrescentar contexto organizacional, o que sugere que fornecer contexto é uma forma barata de buscar mais qualidade sem perder o ganho de tempo. Como discutido na Seção~\ref{sec:resultados}, a dificuldade percebida e as dúvidas relatadas pelos participantes também não se concentram em quem teve o \textit{template} reprovado na execução funcional, o que indica que o quanto a pessoa achou a tarefa difícil não é um bom indício da qualidade final do artefato.

O segundo objetivo era avaliar a conformidade dos \textit{templates} com boas práticas e padrões de IaC, e aqui o resultado depende da etapa observada. Na validação sintática, o manual tem a maior aprovação (93,3\%), seguido pelo com contexto (88,9\%) e pelo sem contexto (62,2\%). Já na execução funcional com LocalStack essa ordem se inverte: o com contexto tem a maior aprovação (95,6\%), seguido pelo sem contexto (91,1\%) e pelo manual (78,3\%). Isso mostra que o contexto organizacional ajuda sobretudo a IA a seguir convenções de formatação, o tipo de regra que um \textit{prompt} mais específico resolve com facilidade. O manual, por outro lado, tende a escrever código sintaticamente correto, mas tropeça mais na execução real: oito das treze reprovações do cenário manual vêm da mesma causa ambiental, quando o \textit{provider} da AWS tenta usar uma operação que o ambiente de teste não suporta, e isso não diz nada sobre a qualidade do \textit{template} em si. Como também discutido na Seção~\ref{sec:resultados}, a relação entre experiência prévia em IaC e aprovação funcional é fraca nesta amostra, e os piores resultados reprovam pelas mesmas causas ambientais que afetam qualquer nível de experiência. Isso sugere que o conhecimento prévio em Terraform conta menos do que se esperava para produzir um \textit{template} funcional, pelo menos nesta amostra. A conformidade ampla medida pelo Checkov não acompanha essa diferença: os três cenários ficam muito próximos (89,3\% a 89,9\%), porque a maior parte das reprovações recai sobre controles que não fazem parte dos requisitos de nenhum grupo, como ciclo de vida e replicação do S3.

O terceiro objetivo era investigar vulnerabilidades de segurança produzidas pelas diferentes formas de geração. A diferença mais clara aparece na saída irrestrita do Security Group, presente em 93,3\% dos \textit{templates} sem contexto, contra 26,7\% com contexto e 45,0\% no manual. Isso mostra que pedir, de forma genérica, para evitar abertura irrestrita não é suficiente para a IA sem contexto, enquanto uma regra organizacional concreta, como proibir \texttt{0.0.0.0/0} fora da porta 443, aproxima o resultado da IA do manual. No S3, a ausência de chave gerenciada pelo cliente aparece no manual (85,0\%) e no com contexto (70,0\%), mas não no sem contexto. Isso não é uma falha de segurança introduzida pela IA, e sim reflexo de uma escolha de criptografia (SSE-S3) que o manual e o cenário com contexto adotaram e o sem contexto não.

No geral, a IA generativa produz \textit{templates} funcionais em uma fração do tempo humano, e apresenta resultados comparáveis ou superiores ao desenvolvimento manual em parte das métricas avaliadas, sobretudo na execução funcional e na redução de configurações inseguras de rede, quando recebe um contexto organizacional explícito. Esse mesmo contexto é o que mais reduz o risco de exposição de rede. Nesta amostra, os resultados sugerem que a qualidade dos \textit{templates}, manuais ou gerados por IA, esteve mais associada a regras explícitas e verificáveis, como nomes de variáveis, restrições de rede e padrões de formatação, do que à experiência prévia de quem os produziu ou à ausência de orientação adicional dada à IA. Por isso, recomenda-se que a adoção de IA na geração de IaC venha acompanhada de um contexto organizacional explícito e da mesma validação automatizada usada no desenvolvimento manual, desenhada para reconhecer variações legítimas de nome e de formato de \textit{output} sem confundi-las com defeito do recurso avaliado.

Como trabalhos futuros, propõe-se um ciclo totalmente agêntico, em que o retorno dos \textit{pipelines} volta para o modelo corrigir o \textit{template} até passar em todas as verificações. Propõe-se também ampliar os cenários para outros recursos da AWS, como EC2, RDS, Lambda e VPC, incluindo arquiteturas com dependência entre recursos, e expressar o contexto organizacional como políticas verificáveis por código. Por fim, um cenário híbrido, com revisão humana sobre \textit{templates} gerados com contexto, e a repetição do experimento com outros modelos e ferramentas de IaC, como OpenTofu e AWS CloudFormation, ajudariam a confirmar o alcance desses resultados.
